\documentclass{article}
\usepackage{geometry}
\geometry{
	a4paper,
	left=2.5cm,   
	right=2.5cm,   
	top=2.5cm,      
	bottom=2.5cm,  
	headheight=13.6pt
}

% 最简单、最稳定的中文方案（自动适配系统，永不报错）
\usepackage[UTF8, scheme=plain]{ctex}

% 设置英文默认字体为 Times New Roman
\usepackage{fontspec}
\setmainfont{Times New Roman}

\usepackage{amsmath}
\usepackage{amssymb}
\usepackage{booktabs}
\usepackage{caption}
\usepackage{setspace}
\onehalfspacing
\setlength{\parskip}{4pt}

\usepackage{array}
\usepackage{listings}
\usepackage{graphicx}
\usepackage{enumitem}
\usepackage{subcaption}
\usepackage{multirow}
\usepackage{mathrsfs}
\usepackage{bm}
\usepackage{xcolor}
\usepackage{colortbl}

\usepackage{tikz}
\usepackage{tikz-3dplot}
\usetikzlibrary{
	arrows.meta, 
	calc, 
	patterns, 
	decorations.pathmorphing, 
	backgrounds, 
	positioning, 
	fit, 
	shapes.geometric, 
	intersections, 
	through, 
	angles, 
	quotes,
	plotmarks
}

%目录超链接
\usepackage{hyperref}
\hypersetup{
	colorlinks=true,
	linkcolor=black,
	filecolor=magenta,      
	urlcolor=cyan,
}
\usepackage{tocloft}
\setlength{\cftbeforesecskip}{2ex}
\setlength{\cftbeforesubsecskip}{1ex}
\setlength{\cftbeforesubsubsecskip}{1ex}

% 定义代码颜色
\definecolor{codegreen}{rgb}{0,0.6,0}
\definecolor{NavyBlue}{RGB}{0,0,128}
\definecolor{PineGreen}{RGB}{1,121,111}

% 配置 listings 环境
\lstset{
	language=TeX,
	tabsize=4,
	captionpos=b,
	numbers=left,
	numbersep=1em,
	sensitive=true,
	showtabs=false,
	frame=shadowbox,
	breaklines=true,
	keepspaces=true,
	showspaces=false,
	showstringspaces=false,
	breakatwhitespace=false,
	basicstyle=\ttfamily\normalsize,
	keywordstyle=\color{NavyBlue},
	commentstyle=\color{codegreen},
	numberstyle=\color{black},
	stringstyle=\color{PineGreen!90!black},
	rulesepcolor=\color{red!20!green!20!blue!20},
	escapeinside={/@}{@/},
}

\definecolor{LightGreen}{rgb}{0.9, 0.95, 0.9}
\definecolor{LightBlue}{RGB}{229, 238, 255}
\definecolor{LightPink}{RGB}{255, 229, 237}
\definecolor{LineBlue}{RGB}{0, 90, 160}
\definecolor{PinRed}{RGB}{190, 30, 45}

\captionsetup[table]{
	justification=centering,
	labelsep=quad,
	font={bf}
}
\captionsetup[figure]{
	justification=centering,
	labelsep=quad,
	font={bf}
}
\captionsetup[subfigure]{
	font={bf}
}

\begin{document}


\begin{center}
	{\LARGE \textbf{复变函数答案}} \\[0.5cm]
	{\large {\kaishu 石继楷 \quad 编写}} \\[0.3cm]
	{\normalsize 2026年9月21日}
\end{center}

\section{习题一}

\noindent
\begin{minipage}[t]{0.5\textwidth}
	
	\textbf{16.}\quad 证明:\\[4pt]
	\textbf{一、$\partial E\cup E\subseteq E'$}\\[2pt]
	$1^\circ$ $E\subseteq E'$:\\
	设 $z_0\in E$.\\
	$\because E$ 为开集,\\
	$\therefore\exists r>0$ 使 $B(z_0,r)\subseteq E$.\\
	$\forall\varepsilon>0$, 取 $w=z_0+\frac{1}{2}\min(r,\varepsilon)$,\\
	则 $0<|w-z_0|<\min(r,\varepsilon)$,\\
	$\therefore w\in E$ 且 $w\in B(z_0,\varepsilon)\setminus\{z_0\}$,\\
	$\therefore z_0\in E'$, 即 $E\subseteq E'$.\\[4pt]
	$2^\circ$ $\partial E\subseteq E'$:\\
	设 $z_0\in\partial E$.\\
	若 $z_0\in E$, 由 $1^\circ$ 知 $z_0\in E'$;\\
	若 $z_0\notin E$,\\
	$\because z_0\in\partial E$,\\
	$\therefore\forall\varepsilon>0$, $B(z_0,\varepsilon)\cap E\neq\varnothing$,\\
	且其中各点均异于 $z_0$,\\
	$\therefore\big(B(z_0,\varepsilon)\setminus\{z_0\}\big)\cap E\neq\varnothing$,\\
	$\therefore z_0\in E'$.\\
	综上, $\partial E\cup E\subseteq E'$.\\[4pt]
	\textbf{二、$E'\subseteq\partial E\cup E$}\\[2pt]
	设 $z_0\in E'$.\\
	若 $z_0\in E$, 结论已成立;\\
	若 $z_0\notin E$,\\
	$\because z_0\in E'$,\\
	$\therefore\forall\varepsilon>0$, $\big(B(z_0,\varepsilon)\setminus\{z_0\}\big)\cap E\neq\varnothing$,\\
	$\therefore B(z_0,\varepsilon)\cap E\neq\varnothing$;\\
	又 $\because z_0\notin E$,\\
	$\therefore z_0\in B(z_0,\varepsilon)\cap E^c$,\\
	$\therefore B(z_0,\varepsilon)\cap E^c\neq\varnothing$,\\
	$\therefore z_0\in\partial E$.\\
	综上, $E'\subseteq\partial E\cup E$.\\[4pt]
	综合一、二, $E'=\partial E\cup E$. $\blacksquare$
	
\end{minipage}
\hfill
\begin{minipage}[t]{0.5\textwidth}
	
	\centering
	\begin{tikzpicture}[baseline=(current bounding box.north), scale=1.5]
		
		% 开圆盘: 虚线边界不属于 E
		\fill[LightBlue,opacity=0.6] (0,0) circle (1.6cm);
		\draw[LineBlue,dashed,line width=0.9pt] (0,0) circle (1.6cm);
		\node[LineBlue,font=\scriptsize] at (-0.75,1.0) {$E$};
		
		% 内点及其邻域: 整个含于 E
		\coordinate (z0) at (0.2,0.0);
		\fill[LineBlue] (z0) circle (1.5pt);
		\draw[LineBlue,line width=0.7pt] (z0) circle (0.45cm);
		\node[LineBlue,font=\scriptsize,anchor=north east] at (-0.05,-0.42) {$z_0\in E$};
		\node[LineBlue,font=\scriptsize,anchor=south west] at (0.52,0.40) {$B(z_0,r)$};
		
		% 边界点及其邻域: 与 E 的交非空(粉色), 与 E^c 的交也非空
		\coordinate (z1) at (45:1.6);
		\begin{scope}
			\clip (0,0) circle (1.6cm);
			\fill[LightPink,opacity=0.9] (z1) circle (0.45cm);
		\end{scope}
		\fill[PinRed] (z1) circle (1.5pt);
		\draw[PinRed,line width=0.7pt] (z1) circle (0.45cm);
		\node[PinRed,font=\scriptsize,anchor=south west] at (1.28,1.34) {$z_1\in\partial E$};
		\node[PinRed,font=\scriptsize,anchor=west] at (1.62,0.95) {$B(z_1,\varepsilon)$};
		
		% 图注
		\node[font=\scriptsize,align=center] at (0,-2.2) {粉色区域为 $B(z_1,\varepsilon)\cap E\neq\varnothing$;\\ 而 $B(z_0,r)\subseteq E$};
		
	\end{tikzpicture}
	
\end{minipage}



\end{document}